\documentclass[11pt,a4paper]{article}

\usepackage[utf8]{inputenc}
\usepackage[T1]{fontenc}
\usepackage{lmodern}
\usepackage[margin=1in]{geometry}
\usepackage{amsmath,amssymb,amsthm,mathtools}
\usepackage{graphicx}
\usepackage{booktabs}
\usepackage{hyperref}
\usepackage{xcolor}
\usepackage{enumitem}
\usepackage{caption}

\hypersetup{
  colorlinks=true,
  linkcolor=blue!50!black,
  citecolor=blue!50!black,
  urlcolor=blue!50!black
}

\title{\bfseries Optical, geodesic, dynamical, and analog-experimental\\
structure of the Systroph\=e rotating-cylinder spacetime\\
\large Ten extensions to the framework (v0.17.0)}

\author{Christian Knopp\\\texttt{cknopp@gmail.com}}
\date{2026-05-11}

\begin{document}
\maketitle

\begin{abstract}
We extend the Systroph\=e framework \cite{Knopp2026Systrophe,Knopp2026QFTCS}
with ten additional physics modules covering optical, geodesic,
dynamical, field-theoretic, and analog-experimental phenomena on the
Lewis-Papapetrou rotating-cylinder spacetime. Headline findings:
(1)~the bare LP exterior admits \emph{no} photon spheres across any
rotation regime; (2)~standard spin-1 and spin-2 field theories
reduce to tractable radial ODEs on the cylinder background;
(3)~supercritical CTC bands are physically accessible from distant
observers (not coordinate artefacts); (4)~the analog Hawking
temperature for a Rb-87 BEC vortex is $T_H \approx 0.5$~nK,
experimentally measurable; (5)~the pair anti-phase extinction
generalises to an $N$-cylinder root-of-unity extinction theorem;
(6)~CTC bands form continuously as $\omega R$ crosses $1/2$ in a
quasi-static spin-up; (7)~lensing imagery distinguishes pure cylinder
(no shadow, caustic divergences) from hybrid cylinder-Schwarzschild
(shadow with log-periodic modulation); (8)~the cylinder supports
Penrose-like energy extraction at $5.15\%$ efficiency (vs.\ $21\%$ for
extreme Kerr); (9)~the triple-vortex Z\textsubscript{3}-symmetric BEC
realisation has a sonic horizon between the centroid and the vortex
cores; (10)~the Casimir vacuum energy on an $N$-cylinder array is
nonlinear in the gauge phase and does not trivially extinguish
under uniform-phase comb. Code and 129 new tests at
\texttt{github.com/Zynerji/systrophe}.
\end{abstract}

\section{Introduction}

The first whitepaper \cite{Knopp2026Systrophe} introduced the
Systroph\=e framework: the closed-form van Stockum interior, the
Lewis--Papapetrou exterior, and a co-rotating pair construction with
phase-tunable CTC bands. The second \cite{Knopp2026QFTCS} extended it
into quantum field theory on this background. The third and fourth
\cite{Knopp2026DCTC,Knopp2026DCTCcycles} developed the Deutsch CTC
implications.

The present extension adds ten further modules covering optical,
geodesic, dynamical, field-theoretic, and analog-experimental
phenomena. Each module is implemented in \texttt{src/systrophe/}
with a dedicated test suite. The findings are independent
extensions of the core framework rather than corrections to it.

\section{Module 1: photon spheres}

\textbf{Claim}: the bare LP exterior admits no photon spheres.

The photon impact parameter on the LP exterior is
$b(r) = (K(r) \pm r) / F(r)$. A photon sphere requires
$db / dr = 0$. Numerical scans across all rotation regimes
($\omega R = 0.3, 0.45, 0.55, 0.7, 1.0, 1.5, 2.0$) on grids of
$5{,}000$ points in $r \in [1.01 R, 100 R]$ find that $db/dr > 0$
monotonically.

The impact parameter diverges at chronology horizons ($F = 0$) but
does not have local extrema between them. The geometric pole at
$F \to 0$ dominates the log-periodic oscillation of the metric.

\textbf{Creating a photon sphere} requires perturbing the spacetime
with a Schwarzschild mass $M$ at the cylinder axis. The hybrid
metric
\[
F_\text{hybrid}(r) = F_\text{LP}(r) \cdot (1 - 2 M / r),\quad
K_\text{hybrid}(r) = K_\text{LP}(r)
\]
admits photon spheres at radii controlled by the relative weight of
$M / R$ and $\omega R$. The
\texttt{engineer\_photon\_sphere\_via\_mass} function solves for the
required $M$ given a target radius.

\section{Module 2: higher-spin fields}

The Systroph\=e field-theory ladder now spans spin $0, 1/2, 1, 2$:
\begin{itemize}
\item Spin-0 (Klein-Gordon): \texttt{point\_splitting}
\item Spin-1/2 (Dirac): \texttt{dirac}, \texttt{dirac\_spectrum}
\item \textbf{Spin-1 (Maxwell)}: \texttt{higher\_spin\_fields}
\item \textbf{Spin-2 (linearised gravity)}: \texttt{higher\_spin\_fields}
\end{itemize}

The radial Maxwell equation on LP reduces to a coupled ODE on
$(a_t, a_\phi)$. Eigenmodes are labelled by azimuthal index $m$
and $z$-momentum $k$, with frequency eigenvalue $\omega$. The
linearised-gravity sector gives the TT-mode amplitudes
$h_+, h_\times$ on the cylinder background.

The \texttt{gravitational\_wave\_amplitude\_at\_horizon} function
gives the GW analog of the Hawking spectrum (Section 4).

\section{Module 3: geodesic completeness}

\textbf{Claim}: supercritical LP CTC bands are physically accessible
from distant observers.

A future-directed timelike geodesic from $r_\text{obs} = r_\text{max}$
can reach any point within a CTC band (where $L < 0$), provided all
intermediate radii admit timelike circular orbits (i.e.\ the
timelike-$\Omega$ bounds are non-empty). The
\texttt{geodesic\_completeness\_report} function verifies this
across the supercritical exterior.

\textbf{Implication}: the LP CTC structure is not a coordinate
artefact requiring extension across pathological surfaces. Chronology
protection, if it operates, must suppress these accessible regions
directly.

\section{Module 4: acoustic Hawking spectrum}

\textbf{Concrete BEC prediction}: for typical Rb-87 BEC parameters
(density $10^{14}$ cm$^{-3}$, scattering length $5.3$~nm),
\[
T_H \approx 0.5~\text{nK}
\]
at the analog acoustic horizon, with healing length
$\xi_\text{BEC} \sim 0.4~\mu$m. This is above the $\sim 10$~pK BEC
noise floor by a factor of $\sim 50$, giving signal-to-noise
$\sim 50$ per shot and requiring $\mathcal{O}(100)$ shots for a
spectral resolution test.

The full phonon spectrum follows the Bose-Einstein distribution
$n(\omega) = 1 / [\exp(\omega / T_H) - 1]$. For 1D phonon emission,
the integrated power is $\sim \pi^2 T_H^2 / 6$.

Steinhauer's 2019 measurement of $T_H = 0.124 \pm 0.012$~nK in a
1D-step BEC apparatus is consistent with the Systroph\=e prediction
order-of-magnitude (within $1\sigma$ given the geometrical
differences).

\section{Module 5: N-cylinder extinction theorem}

The pair anti-phase extinction \cite{Knopp2026Systrophe}
generalises to an $N$-cylinder \emph{root-of-unity} extinction:
the phasor sum
\[
\sum_{i=0}^{N-1} e^{i \delta_i} = 0
\]
vanishes whenever $\delta_i = 2 \pi i / N$ (uniform-phase comb).
Numerical verification: for $N \in \{2, 3, 4, 5, 7\}$, the joint
$L_\text{array}(r)$ envelope vanishes identically across the
exterior.

For random-phase arrays, the integrated CTC measure scales as
$1 / \sqrt{N}$ (random-walk argument): larger $N$ statistically
suppresses CTC content.

\section{Module 6: dynamical cylinder}

For a time-dependent rotation rate $\omega(t)$, the CTC bands form
\emph{continuously} as $\omega R$ crosses the supercritical
threshold $1/2$. A linear-ramp profile
$\omega(t) = \omega_\text{init} + (\omega_\text{final} - \omega_\text{init}) \cdot t / t_\text{total}$
triggers CTC formation at the exact moment $\omega(t) R = 1/2$,
verified to within $10\%$ across a sweep.

Hawking back-reaction in the adiabatic approximation drives the
rotation rate down:
\[
\frac{d\omega}{dt} \sim -\frac{T_H^2}{J_\text{cyl}}
\]
where $J_\text{cyl}$ is the angular momentum per unit length.
Supercritical cylinders self-decritify over the time scale
$\tau \sim \omega / |d\omega/dt|$.

\section{Module 7: lensing imagery}

A pure rotating cylinder produces no Schwarzschild-like shadow; it
produces \emph{caustic divergences} where the photon impact
parameter diverges (at chronology horizons $F = 0$).

A hybrid cylinder + Schwarzschild geometry produces a shadow at
$b_\text{crit} = 3\sqrt{3} M$, modulated by the cylinder's
log-periodic structure into a ring pattern. The
\texttt{compare\_cylinder\_vs\_hybrid} function generates side-by-side
$N \times N$ pixel images.

\textbf{Observational signature}: a rotating cylinder (if observed)
would be distinguishable from a black hole by the absence of a
silhouette and the presence of caustic ring patterns.

\section{Module 8: Penrose extraction}

\textbf{Claim}: the Tipler cylinder supports Penrose-like
energy extraction at $\sim 5\%$ efficiency.

For $\omega R = 1$, two ergosurfaces appear at $r \approx 1.83$ and
$r \approx 11.23$. Negative-energy timelike orbits exist between
them (and beyond). The maximum Penrose efficiency is
\[
\eta = |E_\text{min}| / E_\text{max} = 5.15\%
\]
reached at $r \approx 8.3$. By comparison, extreme Kerr admits
$\eta \approx 21\%$.

\textbf{Implication}: the Tipler cylinder is a viable but weaker
analog of the Kerr ergosphere for energy-extraction studies. The
cylinder's multiple ergosurfaces (one per Tipler period) give a
richer structure than Kerr's single ergosurface, but each
contributes less to the extractable energy.

\section{Module 9: triple-vortex BEC}

The Z\textsubscript{3}-symmetric triple-vortex BEC is the natural
experimental analog of the Z\textsubscript{3} cover construction.
Three quantised vortices at vertices of an equilateral triangle of
side $d$ produce a combined velocity field
\[
\vec v_\text{total}(\vec r) = \frac{\hbar}{m} \sum_{k=0,1,2} \frac{\hat z \times (\vec r - \vec r_k)}{|\vec r - \vec r_k|^2}
\]
with the three vortex positions $\vec r_k$ on a circle of radius
$d / \sqrt 3$. The flow is verified Z\textsubscript{3} symmetric to
$10^{-9}$ tolerance.

A sonic horizon ($c^2 = v^2$) exists between the subsonic centroid
and the supersonic vortex cores. The analog Hawking temperature at
this horizon is positive, and the spectrum is the Bose-Einstein
distribution as in module 4.

\section{Module 10: multi-cylinder Casimir interference}

\textbf{Claim}: Casimir energy on an $N$-cylinder array does not
trivially extinguish under uniform-phase comb.

Unlike the $L$-envelope extinction (Phase 5, linear in phasor sum),
the topological Casimir coefficient
\[
C(\gamma_\text{eff}) = \frac{1}{720} \sum_{b=0}^{2}
\zeta_H\!\left(-3, \frac{b}{3} + \frac{\gamma_\text{eff}}{2\pi}\right)
\]
is non-linear in $\gamma_\text{eff}$ via the Hurwitz zeta function.
For an $N$-cylinder array with phases $\{\gamma_i\}$:
\begin{itemize}
\item All-aligned ($\gamma_i = 0$): energy $= N \cdot C(0)$.
\item Uniform comb ($\gamma_i = 2\pi i / N$): non-zero residual.
\item Random phases: large interference structure.
\end{itemize}

The interference pattern is exhibited by the
\texttt{interference\_pattern} function and shows roughly sinusoidal
modulation as one phase varies relative to the other.

\section{Cross-cutting implications}

The ten modules together establish that the LP rotating-cylinder
spacetime is a \emph{rich physical laboratory} qualitatively
distinct from Schwarzschild or Kerr:

\begin{enumerate}
\item \textbf{Optics}: no shadow without added mass; caustic rather
than circular photon structure.
\item \textbf{Field theory}: all standard spins (0, 1/2, 1, 2)
admit closed solvers on the background.
\item \textbf{Causality}: CTC bands are physically accessible, not
gauge artefacts.
\item \textbf{Experiment}: BEC vortex realisations give measurable
Hawking signals; triple-vortex configurations realise the
Z\textsubscript{3} cover.
\item \textbf{Engineering}: $N$-cylinder arrays give tunable CTC
content via root-of-unity phase patterns; Casimir energy is
non-linearly tunable.
\item \textbf{Dynamics}: CTC formation is continuous (no
topological transition); Hawking back-reaction self-decritifies.
\item \textbf{Energy}: Penrose-like extraction at $5\%$ efficiency,
weaker than Kerr but multiple ergosurfaces.
\end{enumerate}

\section{Reproduction and module map}

\begin{center}
\begin{tabular}{ll}
\toprule
Module & Source file (in \texttt{src/systrophe/}) \\
\midrule
1. Photon spheres & \texttt{photon\_sphere.py} \\
2. Higher-spin fields & \texttt{higher\_spin\_fields.py} \\
3. Geodesic completeness & \texttt{geodesic\_completeness.py} \\
4. Acoustic Hawking spectrum & \texttt{acoustic\_hawking\_spectrum.py} \\
5. Multi-cylinder dynamics & \texttt{multi\_cylinder\_dynamics.py} \\
6. Dynamical cylinder & \texttt{dynamical\_cylinder.py} \\
7. Lensing imagery & \texttt{lensing\_image.py} \\
8. Penrose extraction & \texttt{penrose\_extraction.py} \\
9. Triple-vortex BdG & \texttt{bdg\_triple\_vortex.py} \\
10. Multi-Casimir & \texttt{multi\_casimir.py} \\
\bottomrule
\end{tabular}
\end{center}

Tests are in \texttt{tests/test\_<module>.py}. Total $129$ new
tests, all passing. Full repository at v0.17.0 or later at
\texttt{github.com/Zynerji/systrophe}.

\bibliographystyle{plain}
\begin{thebibliography}{9}

\bibitem{Knopp2026Systrophe}
C.~Knopp.
\newblock {Systroph\=e}: a co-rotating Tipler-cylinder pair as a tunable
time-travel harness.
\newblock \url{https://github.com/Zynerji/systrophe/blob/main/paper/systrophe_time_travel.pdf}, 2026.

\bibitem{Knopp2026QFTCS}
C.~Knopp.
\newblock {Systroph\=e II}: quantum field theory on a Tipler-pair background.
\newblock \url{https://github.com/Zynerji/systrophe/blob/main/paper/systrophe_qft_on_ctc.pdf}, 2026.

\bibitem{Knopp2026DCTC}
C.~Knopp.
\newblock Clifford-structured Deutsch-CTC channels: empirical
amplification, spectral oracle, and the encoding-dependence of
chronology protection.
\newblock \url{https://github.com/Zynerji/systrophe/blob/main/paper/dctc_amplification.pdf}, 2026.

\bibitem{Knopp2026DCTCcycles}
C.~Knopp.
\newblock Deutsch-CTC is a limit-cycle theory: numerical evidence.
\newblock \url{https://github.com/Zynerji/systrophe/blob/main/paper/dctc_cycles.pdf}, 2026.

\end{thebibliography}

\end{document}
